\documentclass[UTF8,fontset=fandol,11pt]{ctexart}
\usepackage[a4paper,margin=22mm]{geometry}
\usepackage{amsmath,amssymb,bm,mathtools}
\newcommand{\ket}[1]{\left|#1\right\rangle}
\newcommand{\bra}[1]{\left\langle#1\right|}
\newcommand{\braket}[1]{\left\langle#1\right\rangle}
\DeclareMathOperator{\Tr}{Tr}
\DeclareMathOperator{\tr}{tr}
\setlength{\parindent}{0pt}
\setlength{\parskip}{0.5em}
\pagestyle{plain}
\begin{document}
\section*{译文 · 简体中文}
对于参数 \(\theta\)，量子态为 \(\rho_\theta\)。量子 Fisher 信息满足
\par
\[F_Q(\rho_\theta)=\operatorname{Tr}(\rho_\theta L_\theta^2),\quad \partial_\theta \rho_\theta=\frac{L_\theta\rho_\theta+\rho_\theta L_\theta}{2}.\]
\par
原文给出的概率满足 50\% 的示例条件，保留符号 a\_b \& c。
\end{document}
